% GENERATED FILE. Edit canonical lessons, then run npm run generate:books.
% canonical-source-hash: fnv1a64:7d619435577ef847
% canonical-lessons: IT-C504-mazzo, IT-C504-bluffare, IT-C504-scommettere, IT-C504-scommessa, IT-C504-cruciverba

\chapter{Carte e scommesse}
\label{ch:it-carte-e-scommesse}
\hlchaptermodality{504}

\begin{chapteropening}
\textbf{By the end of this chapter:} \emph{I can talk about card games, bets and crosswords.}

Complete the last lesson of chapter 504: I can talk about card games, bets and crosswords.
\end{chapteropening}

\section[il mazzo]{il mazzo — deck (of cards); bunch}

\label{lesson:IT-C504-mazzo}



\begin{quote}
Before the new one: say the Italian for board game, then the Italian for to cheat.
\end{quote}

\subsection*{You'll want to know: il mazzo}
\textbf{il mazzo} — \textquotedblleft{}deck (of cards); bunch\textquotedblright{}.

\begin{quote}
Il nonno mescola il mazzo e dà le carte.
\end{quote}

\begin{quote}
\emph{Grandad shuffles the deck and deals the cards.}
\end{quote}

\emph{Un mazzo di fiori} — a bunch of flowers; \emph{un mazzo di chiavi} — a bunch of keys.

\begin{cousinweb}
\emph{mazzo} is usually traced, like \emph{mazza} \textquotedblleft{}a club\textquotedblright{}, to Vulgar Latin \emph{mattea}, from Latin \emph{mateola} \textquotedblleft{}a mallet\textquotedblright{} (EtymDB).
\end{cousinweb}

\begin{grammarlens}[title={What you've built}]
One more word for telling, explaining and describing at B1.
\end{grammarlens}

\subsection*{Your turn}
\begin{itemize}
\raggedright
  \item \emph{Say it:} \emph{il mazzo}
  \item \emph{Say it:} \emph{il mazzo}, once more
  \item \emph{Say it:} a sentence of your own with \emph{il mazzo}
  \item \emph{Recall:} say the Italian for board game, then the Italian for to cheat, then say \emph{il mazzo} again
\end{itemize}

\begin{tcolorbox}[breakable,colback=teal!4,colframe=teal!35!black,title={Before you move on}]
What does il mazzo mean? (\textquotedblleft{}Deck (of cards); bunch\textquotedblright{}.) Say it once more.
\end{tcolorbox}
\section[bluffare]{bluffare — to bluff}

\label{lesson:IT-C504-bluffare}



\begin{quote}
Before the new one: say the Italian for to cheat, then the Italian for deck.
\end{quote}

\subsection*{You'll want to know: bluffare}
\textbf{bluffare} — \textquotedblleft{}to bluff\textquotedblright{}.

\begin{quote}
Non ha niente in mano, sta bluffando.
\end{quote}

\begin{quote}
\emph{He has nothing in his hand, he is bluffing.}
\end{quote}

Also \emph{fare un bluff}.

\begin{cousinweb}
\emph{bluffare} is built on \emph{bluff}, which was borrowed from English \emph{bluff} (EtymDB), a word from card games.
\end{cousinweb}

\begin{grammarlens}[title={What you've built}]
One more word for telling, explaining and describing at B1.
\end{grammarlens}

\subsection*{Your turn}
\begin{itemize}
\raggedright
  \item \emph{Say it:} \emph{bluffare}
  \item \emph{Say it:} \emph{bluffare}, once more
  \item \emph{Say it:} a sentence of your own with \emph{bluffare}
  \item \emph{Recall:} say the Italian for to cheat, then the Italian for deck, then say \emph{bluffare} again
\end{itemize}

\begin{tcolorbox}[breakable,colback=teal!4,colframe=teal!35!black,title={Before you move on}]
What does bluffare mean? (\textquotedblleft{}To bluff\textquotedblright{}.) Say it once more.
\end{tcolorbox}
\section[scommettere]{scommettere — to bet}

\label{lesson:IT-C504-scommettere}



\begin{quote}
Before the new one: say the Italian for deck, then the Italian for to bluff.
\end{quote}

\subsection*{You'll want to know: scommettere}
\textbf{scommettere} — \textquotedblleft{}to bet\textquotedblright{}.

\begin{quote}
Scommetto che domani piove.
\end{quote}

\begin{quote}
\emph{I bet it rains tomorrow.}
\end{quote}

The past participle is \emph{scommesso}. \emph{Scommetto che…} — I bet that…

\begin{cousinweb}
\emph{scommettere} is built from \emph{s-} and \emph{commettere} (EtymDB), inherited from Latin \emph{committere} \textquotedblleft{}to bring together, to match\textquotedblright{}, from \emph{mittere} \textquotedblleft{}to send\textquotedblright{}.
\end{cousinweb}

\begin{grammarlens}[title={What you've built}]
One more word for telling, explaining and describing at B1.
\end{grammarlens}

\subsection*{Your turn}
\begin{itemize}
\raggedright
  \item \emph{Say it:} \emph{scommettere}
  \item \emph{Say it:} \emph{scommettere}, once more
  \item \emph{Say it:} a sentence of your own with \emph{scommettere}
  \item \emph{Recall:} say the Italian for deck, then the Italian for to bluff, then say \emph{scommettere} again
\end{itemize}

\begin{tcolorbox}[breakable,colback=teal!4,colframe=teal!35!black,title={Before you move on}]
What does scommettere mean? (\textquotedblleft{}To bet\textquotedblright{}.) Say it once more.
\end{tcolorbox}
\section[la scommessa]{la scommessa — bet}

\label{lesson:IT-C504-scommessa}



\begin{quote}
Before the new one: say the Italian for to bluff, then the Italian for to bet.
\end{quote}

\subsection*{You'll want to know: la scommessa}
\textbf{la scommessa} — \textquotedblleft{}bet\textquotedblright{}.

\begin{quote}
Ho vinto la scommessa: mi devi un caffè.
\end{quote}

\begin{quote}
\emph{I won the bet: you owe me a coffee.}
\end{quote}

\emph{Fare una scommessa} — to make a bet.

\begin{cousinweb}
\emph{scommessa} is the feminine past participle of \emph{scommettere} \textquotedblleft{}to bet\textquotedblright{}, turned into a noun; the verb goes back to Latin \emph{committere} (EtymDB).
\end{cousinweb}

\begin{grammarlens}[title={What you've built}]
One more word for telling, explaining and describing at B1.
\end{grammarlens}

\subsection*{Your turn}
\begin{itemize}
\raggedright
  \item \emph{Say it:} \emph{la scommessa}
  \item \emph{Say it:} \emph{la scommessa}, once more
  \item \emph{Say it:} a sentence of your own with \emph{la scommessa}
  \item \emph{Recall:} say the Italian for to bluff, then the Italian for to bet, then say \emph{la scommessa} again
\end{itemize}

\begin{tcolorbox}[breakable,colback=teal!4,colframe=teal!35!black,title={Before you move on}]
What does la scommessa mean? (\textquotedblleft{}Bet\textquotedblright{}.) Say it once more.
\end{tcolorbox}
\section[il cruciverba]{il cruciverba — crossword}

\label{lesson:IT-C504-cruciverba}



\begin{quote}
Before the new one: say the Italian for to bet, then the Italian for bet.
\end{quote}

\subsection*{You'll want to know: il cruciverba}
\textbf{il cruciverba} — \textquotedblleft{}crossword\textquotedblright{}.

\begin{quote}
In treno faccio sempre il cruciverba.
\end{quote}

\begin{quote}
\emph{On the train I always do the crossword.}
\end{quote}

The same form in the plural: \emph{i cruciverba}.

\begin{cousinweb}
\emph{cruciverba} is a modern word built from Latin \emph{crux, crucis} \textquotedblleft{}a cross\textquotedblright{} and \emph{verba} \textquotedblleft{}words\textquotedblright{} (Lewis \& Short).
\end{cousinweb}

\begin{grammarlens}[title={What you've built}]
One more word for telling, explaining and describing at B1.
\end{grammarlens}

\subsection*{Your turn}
\begin{itemize}
\raggedright
  \item \emph{Say it:} \emph{il cruciverba}
  \item \emph{Say it:} \emph{il cruciverba}, once more
  \item \emph{Say it:} a sentence of your own with \emph{il cruciverba}
  \item \emph{Recall:} say the Italian for to bet, then the Italian for bet, then say \emph{il cruciverba} again
\end{itemize}

\begin{tcolorbox}[breakable,colback=teal!4,colframe=teal!35!black,title={Before you move on}]
What does il cruciverba mean? (\textquotedblleft{}Crossword\textquotedblright{}.) Say it once more.
\end{tcolorbox}
